\chapter{結論と限界}\label{chap:7}

本章では、本報告における設計統合と検証から得られた知見を総括し、現在の検証範囲および未検証事項を整理する。
最後に、再現性を確保するための検証対象バージョンと根拠資料の対応を示す。

\section{検証で得られた知見}\label{sec:7-1}

本報告では、既存の2コアCPUに入力ROMのデータ読み出し経路とコアIDの互換アドレスを統合し、共有メモリ上での分割バブルソートの動作とサイクル数をXSim 2023.1で検証した。
選定した18種類の入力ケース（計54条件）すべてにおいて、共有データポートを調停しながらCore0・Core1が各担当区間を並行してソートし、両側の完了同期を経てCore0が正しくマージを完了することを確認した。
さらに、10要素の代表入力に関する波形解析により、競合発生時のCore1のストールとPC保持、待機解除後の実行再開、完了フラグの推移、および最終配列出力に至る一連の挙動を実コードと対応付けて実証した。

サイクル数の比較分析では、同一アルゴリズムを1コアで逐次実行する場合（serial\_split）に比べ、2コア実行（dual）は指定10要素で442から366サイクルへ、固定乱数64要素で8,614から5,527サイクル（35.8\%減）へと短縮された。
処理区間ごとの内訳から、短縮効果の大半はREADYから両側ソート完了までの並行区間に集中している一方、初期コピーおよびマージ・書き戻しを含む逐次区間はほぼ変わらず、共有ポート競合とあわせて2倍速に達しない主要因となっていることが明らかになった。

\section{検証範囲と未検証事項}\label{sec:7-2}

本検証にはいくつかの前提と限界がある。
機能検証の対象は、リセットおよびメモリ初期化後の状態から起動し、選定した最大64要素の符号なし32ビット入力を処理する条件に限られる。
したがって、前回実行のフラグが残存した状態からの再起動、64要素を超える大規模入力、Core0固定優先調停のもとでCore1が長時間待たされる状況、複数の乱数シードによるサイクル数のばらつき、ならびに検査コードのカバレッジ評価は今後の検証課題として残されている。

また、本報告の性能評価はすべてRTLビヘイビアシミュレーション上のサイクル数に基づくものであり、論理合成、配置配線、およびFPGA実機での動作測定は行っていない。
実機における処理時間の短縮率はサイクル数の比だけでなく最大動作周波数にも依存するため、非同期読み出しメモリのFPGA資源への実装方法、回路規模、およびタイミング制約を含めた実機評価が別途必要である。

\section{検証対象の版と根拠資料}\label{sec:7-3}

本報告の検証対象は、原本リポジトリtakehaya/antikytheraの未コミット変更を含む作業状態（HEAD:\ \allowbreak 8b61da20522f5bfcb01908f2342abfea7676bb25）を出発点とし、第\ref{chap:2}章の統合変更を加えた版である。
そのため、HEADコミット単体からは同一の検証環境を再現できない。
参照した全19ファイルの一覧とSHA-256ハッシュ値はsource-manifest.jsonに記録しており、ファイル集合識別子の先頭16桁は62eb51779b0c0173である。

表\ref{tab:8}に本文の記述と保存資料の対応をまとめる。
原本リポジトリは文献\cite{antikythera}に、Vivadoシミュレーションおよび波形操作の参照資料は文献\cite{ug900,ug835}に示す。
なお、参照した公式マニュアルは2020.2版であり、本報告の実測に用いた実行環境（XSim 2023.1）とは版が異なる。

\begin{table}[H]
\centering
\caption{本文の根拠と保存資料の対応}
\label{tab:8}
\footnotesize
\setlength{\tabcolsep}{3pt}
\renewcommand{\arraystretch}{1.25}
\begin{tabular}{>{\raggedright\arraybackslash}p{\dimexpr0.340493\linewidth-2\tabcolsep\relax}>{\raggedright\arraybackslash}p{\dimexpr0.659507\linewidth-2\tabcolsep\relax}}\toprule
\textbf{根拠項目} & \textbf{保存資料}\\\midrule
実装と命令一覧 & work/内のRTL、tests/sort/build\_programs.py、research/cases/original\_ten/dual/listing.s\\
54条件の数値・判定 & report/xsim-results.csv、runs/matrix-*/simulate.logおよび終了記録\\
図\ref{fig:3}〜\ref{fig:5}の波形画面 & runs/ten-dual-v2/evidence.wdb、report/gui-captures-20261011/内の原本と切出し記録\\
途中状態の検査コード & tests/sort/tb\_sort.v、research/sim/tb\_vivado\_sort.v\\
\bottomrule\end{tabular}
\end{table}
